\documentclass[a4paper,11pt]{article}
\usepackage{mysty}
\usepackage[margin=24mm]{geometry}
\usepackage{booktabs}
\usepackage{algorithm}
\usepackage{algpseudocode}
\usepackage[hidelinks,unicode]{hyperref}

\hypersetup{
  pdftitle={信念集合で条件付けた連続行動拡散方策とDPPO型PPO目的},
  pdfauthor={Belief-Score-Based-Model}
}
\floatname{algorithm}{アルゴリズム}
\algrenewcommand\algorithmicrequire{\textbf{入力:}}
\algrenewcommand\algorithmicensure{\textbf{出力:}}
\algrenewcommand\algorithmicfor{\textbf{for}}
\algrenewcommand\algorithmicdo{\textbf{do}}
\algrenewcommand\algorithmicend{\textbf{end}}
\algrenewcommand\algorithmicif{\textbf{if}}
\algrenewcommand\algorithmicthen{\textbf{then}}
\newcommand{\stopgrad}{\operatorname{stopgrad}}
\newcommand{\clip}{\operatorname{clip}}
\newcommand{\Normal}{\mathcal{N}}
\newcommand{\old}{\mathrm{old}}
\newcommand{\env}{\mathrm{env}}
\newcommand{\den}{\mathrm{den}}
\setlength{\emergencystretch}{3em}

\title{信念集合で条件付けた連続行動拡散方策と\\DPPO型PPO目的}
\author{Belief-Score-Based-Model 実験仕様}
\date{\today}

\begin{document}
\maketitle

\begin{abstract}
本稿は、スコアで表した粒子信念を条件として連続行動を生成するDDPM方策と、
その各逆拡散遷移へPPOを適用する学習目的を定義する。
最終行動の周辺尤度は一般に計算不能であるため、DPPO\cite{ren2025dppo}と同様に
逆拡散過程を内側のMDPとして展開し、各Gaussian遷移の尤度比を個別にclipする。
一方、本実験はデモンストレーションによる事前学習を用いないfrom-scratch設定であり、
標準DPPOの主設定とは異なる。信念energyは、rollout時に保存したGaussian noiseから
現在parametersでULAとbelief recurrenceを再parameterizeしてend-to-endに更新する。
本稿ではepisode全体へ逆伝播する\texttt{full}と、環境時刻境界だけを切る
\texttt{tbptt\_1}を区別して実装対象を明示する。active productionでは、
score-belief TransformerとGRUの双方へ同一のDiffusionPolicyとPPO目的を接続し、
action familyを揃えて表現・記憶機構だけを比較する。
\end{abstract}

\section{記法と信念条件}

環境時刻を$t$、観測を$o_t$、環境報酬を$r_t$、連続行動を
$a_t\in\mathcal A\subset\mathbb R^{d_a}$とする。
信念は$K_b$個の粒子と、その位置で評価したscoreの集合
\begin{equation}
  B_t=\left\{\left(s_{t,i},q_{t,i}\right)\right\}_{i=1}^{K_b},
  \qquad q_{t,i}=\mathfrak{s}_t(s_{t,i})\in\mathbb R^{d_s}
  \label{eq:belief-set}
\end{equation}
で近似する。粒子順序に依存しないSet Transformer encoderとpoolingを用いて
\begin{equation}
  c_t=E_\phi(B_t)
  =\operatorname{Pool}\!\left(
    \operatorname{SetTransformer}_\phi
    \bigl([s_{t,i},q_{t,i}]_{i=1}^{K_b}\bigr)
  \right)\in\mathbb R^{d_c}
  \label{eq:belief-encoding}
\end{equation}
を作る。criticと拡散方策はともに$c_t$で条件付ける。
単一行動を生成する場合を記述するが、$a_t$を
$H_a\times d_a$次元のaction chunkに置き換えても数式は同じである。

active比較のGRUでは、部分観測と正規化した直前行動を入力して
\begin{equation}
  h_t=\operatorname{GRU}_{\omega}([o_t,a_{t-1}],h_{t-1}),
  \qquad c_t=W_ch_t
  \label{eq:gru-condition}
\end{equation}
を作る。式\eqref{eq:belief-encoding}と式\eqref{eq:gru-condition}の後段には、
同一構成のcriticとDiffusionPolicyを接続する。reverse step数、noise schedule、
variance floor、denoising advantage重み、per-step PPO clippingも共通である。
したがって連続taskのactive比較はdiffusion対Gaussianではなく、
score-belief表現対GRU記憶の比較である。離散CartPoleでは両者に同じcategorical headを使う。

\section{信念条件付きDDPM方策}

\subsection{forward processとreverse process}

拡散step数を$N$とし、$\beta_n\in(0,1)$、
$\alpha_n=1-\beta_n$、$\bar\alpha_n=\prod_{j=1}^{n}\alpha_j$と置く。
正規化済み行動$x^0$に対するforward processは
\begin{align}
  q(x^n\mid x^{n-1})
    &=\Normal\!\left(\sqrt{\alpha_n}x^{n-1},\beta_n I\right),\\
  q(x^n\mid x^0)
    &=\Normal\!\left(\sqrt{\bar\alpha_n}x^0,
      (1-\bar\alpha_n)I\right)
  \label{eq:forward-ddpm}
\end{align}
である。方策として使うreverse chainは
\begin{equation}
  x_t^N\sim\Normal(0,I),\qquad
  p_\theta(x_t^{n-1}\mid x_t^n,c_t)
  =\Normal\!\left(\mu_\theta(x_t^n,n,c_t),\sigma_n^2 I\right),
  \quad n=N,\ldots,1,
  \label{eq:reverse-kernel}
\end{equation}
とする。noise-prediction parameterizationでは
\begin{equation}
  \mu_\theta(x^n,n,c)
  =\frac{1}{\sqrt{\alpha_n}}
    \left(
      x^n-\frac{\beta_n}{\sqrt{1-\bar\alpha_n}}
      \epsilon_\theta(x^n,n,c)
    \right).
  \label{eq:ddpm-mean}
\end{equation}
実装ではposterior varianceに下限$\sigma_{\min}^2>0$を設ける。
これは最終stepを含む各遷移のGaussian log-densityを有限に保ち、
学習中の探索を確保するためである。

環境へ渡す有界行動は、既知の決定論的変換
\begin{equation}
  a_t=T(x_t^0)
  =\tanh(x_t^0)\odot\frac{a_{\max}-a_{\min}}{2}
   +\frac{a_{\max}+a_{\min}}{2}
  \label{eq:action-transform}
\end{equation}
で作る。PPOで評価する確率変数は、変換後の周辺分布ではなく
式\eqref{eq:reverse-kernel}の各内部遷移である。

\subsection{拡散過程を内側MDPとして扱う}

環境時刻$t$ごとに、$n=N,N-1,\ldots,1$の$N$個の内側遷移を置く。
内側状態、内側行動、報酬をそれぞれ
\begin{equation}
  \bar s_{t,n}=(c_t,x_t^n),\qquad
  \bar a_{t,n}=x_t^{n-1},\qquad
  \bar r_{t,n}=
  \begin{cases}
    0,&n>1,\\
    r_t,&n=1
  \end{cases}
  \label{eq:augmented-mdp}
\end{equation}
とする。したがって各内側方策
$p_\theta(x_t^{n-1}\mid x_t^n,c_t)$は明示的なGaussian likelihoodを持つ。
最終行動の周辺密度
\begin{equation}
  \pi_\theta(x_t^0\mid c_t)
  =\int p(x_t^N)\prod_{n=1}^{N}
    p_\theta(x_t^{n-1}\mid x_t^n,c_t)\,
    \mathrm{d}x_t^{1:N}
  \label{eq:intractable-marginal}
\end{equation}
は計算せず、trajectory上で観測された各遷移をPPO sampleとして使う。

\section{DPPO型clipped PPO目的}

\subsection{各denoising transitionの尤度比}

rollout時のparametersを$(\theta_{\old},\phi_{\old})$とする。
保存したreverse chain $(x_t^n,x_t^{n-1})$とold log-density、および
式\eqref{eq:belief-replay}でcurrent parametersから再構成する$c_t(\phi)$に対し、
各stepのlog ratioを
\begin{align}
  \ell_{t,n}(\theta,\phi)
    &=\log p_\theta(x_t^{n-1}\mid x_t^n,c_t(\phi)),\\
  \rho_{t,n}(\theta,\phi)
    &=\exp\!\left(
      \ell_{t,n}(\theta,\phi)-\ell_{t,n}(\theta_{\old},\phi_{\old})
    \right)
  \label{eq:step-ratio}
\end{align}
と定める。重要なのは、\textbf{各stepのratioを個別にclipする}ことである。
式\eqref{eq:intractable-marginal}のratioも、
$\prod_n\rho_{t,n}$もPPO目的には使用しない。

\subsection{環境GAE、critic、denoising advantage}

criticは内側のnoisy actionには依存させず、環境時刻ごとに
\begin{equation}
  V_\psi(c_t)\simeq
  \mathbb E\!\left[\sum_{j\geq0}\gamma_{\env}^j r_{t+j}\mid c_t\right]
  \label{eq:critic}
\end{equation}
を推定する。episode終了indicatorを$d_t\in\{0,1\}$として
\begin{align}
  \delta_t
    &=r_t+\gamma_{\env}(1-d_t)V_{\psi_{\old}}(c_{t+1})
      -V_{\psi_{\old}}(c_t),\\
  \widehat A_t^{\env}
    &=\sum_{j=0}^{T-t-1}(\gamma_{\env}\lambda)^j
      \left(\prod_{m=0}^{j-1}(1-d_{t+m})\right)\delta_{t+j},\\
  \widehat R_t
    &=\widehat A_t^{\env}+V_{\psi_{\old}}(c_t)
  \label{eq:gae}
\end{align}
とする。$d_t$は既定では真の終端とhorizon打ち切り（truncation）の両方で$1$であり、
打ち切りを終端とみなす規約である。この規約は全手法・全modeで共通なので内部比較の公平性は
保たれるが、絶対returnやvalue学習はtruncation-bootstrap実装と直接比較できない。
設定\texttt{ppo.bootstrap\_on\_truncation}を真にすると、truncationした$t$に限り
$\delta_t$へ$\gamma_{\env}V_{\psi_{\old}}(c_T)$を加える（Pardoら\cite{pardo2018timelimits}の方式）。
$c_T$はauto-resetが捨てる打ち切り直後の真の観測から、実行した行動を条件として
belief（baselineではGRU hidden）を1step進めた条件である。$d_t$自体は変えないので
$\lambda$-returnの累積はepisode境界で切れたままであり、bootstrap項だけが増える。
既定は偽で、投入済みcampaignと同一の挙動になる。この設定はvalue targetそのものを変えるため、
値の異なるrunを同一の集計・比較へ混在させてはならない。

規約が全手法で共通であることは、同一の目的関数を最適化しているという手続き上の公平性を
保証するだけであり、打ち切りbiasの負担が両表現で等しいことまでは意味しない。
式\eqref{eq:gru-condition}の$h_t$は制約のないrecurrent vectorであり、観測と直前行動の
prefixから経過stepを暗黙に数えてhorizon依存のvalueをfitできる。一方、
式\eqref{eq:belief-set}の$B_t$は状態上の事後粒子とそのscoreの集合に限られ、
時間channelを持たない（本実験のどのtaskの部分観測にも時刻は含まれない）。
したがって既定規約下でのvalue学習と絶対returnの解釈には、この非対称性を明示して
織り込む必要がある。内側stepには同じ環境advantageを複製し、
\begin{equation}
  \widehat A_{t,n}^{\den}
  =\gamma_{\den}^{n-1}\widehat A_t^{\env},
  \qquad 0<\gamma_{\den}\leq1
  \label{eq:denoise-advantage}
\end{equation}
と重み付けする。$n=N$が最初で最もnoisyな遷移なので、
$\gamma_{\den}<1$では最初のstepほど強くdownweightされ、$n=1$の最終stepの重みは1になる。

PPOのsurrogateとvalue lossは
\begin{align}
  L_{\mathrm{clip}}(\theta,\phi)
  &=\frac{1}{TN}\sum_{t,n}
    \min\!\left\{
      \rho_{t,n}\widehat A_{t,n}^{\den},
      \clip(\rho_{t,n},1-\varepsilon,1+\varepsilon)
      \widehat A_{t,n}^{\den}
    \right\},\\
  L_V(\psi,\phi)
  &=\frac{1}{2T}\sum_t
    \left(V_\psi(c_t(\phi))-\widehat R_t\right)^2,\\
  L_{\mathrm{total}}
  &=-L_{\mathrm{clip}}+c_VL_V-c_H\mathcal H
  \label{eq:ppo-total}
\end{align}
である。$L_{\mathrm{clip}}$が方策更新項、$L_V$がcriticのvalue regression、
$\mathcal H$が方策entropyである。これらは標準的なPPOの三要素であるが、
本比較実験ではdebug・productionともに$c_H=0$とし、entropy bonusを最適化項から除外する。
activeな連続比較は両手法とも固定scheduleのDiffusionPolicyであり、最適化する目的は
$-L_{\mathrm{clip}}+c_VL_V$で共通である。
実装ではPPOのclipped objectiveと
gradient norm clippingを用いる。approximate KLとclip fractionは診断値として記録するだけで、
更新の打ち切りやminibatchの拒否には使用しない。またlog ratioに独自のhard clampを置かず、
式\eqref{eq:step-ratio}とPPO clipをそのまま用いる。
標準ULAのdriftや状態を変えるscore/particle clampもmain設定では0（無効）とする。
pilotで不安定性が出た場合は、まず学習率とULA step sizeを明示的に再設定する。

\section{再parameterizeしたbelief recurrenceの全微分}

\subsection{明示noiseによる再生}

信念energyのparametersを$\eta=(\eta_0,\eta_f,\eta_u)$とする。初期時刻専用の
energy networkは通常時の観測・遷移energyとはparametersを共有せず、実装するpotentialは
\begin{equation}
  E_\eta(s;o_t,B_{t-1},a_{t-1})=
  \begin{cases}
    g_{\eta_0}(o_0,s)-\frac12\lVert s\rVert^2,
      & \text{episode初期時刻},\\
    f_{\eta_f}(o_t,s,a_{t-1})
      +\operatorname{logmeanexp}_{s^-\in B_{t-1}}
        u_{\eta_u}(s,s^-,a_{t-1}),
      & \text{それ以外}
  \end{cases}
  \label{eq:belief-energy-piecewise}
\end{equation}
とする。
$-\lVert s\rVert^2/2$はparameter regularizationではなく、初期密度の裾とlatent
scaleを固定する標準Gaussian基準密度のlog densityである。
無制限の関数族ならこの項を$g_{\eta_0}$へ吸収する再parameterizationは可能だが、
有限幅MLPが非有界なlatent空間全体で同じ二次tailを正確に表す保証はない。
そのため実装ではnormalizableな裾とlatent scaleを明示する基準項として分離する。
rollout時には粒子座標をPPO更新用の
固定入力として保存するのではなく、各環境時刻$t$についてULA初期noise
$z_{t,i}\sim\Normal(0,I)$と各stepのnoise
$\epsilon_{t,l,i}\sim\Normal(0,I)$を保存する。PPO更新時には観測列、
実現済み行動列、episode-start mask、およびこれらのnoiseを固定し、
現在の$\eta$から
\begin{align}
  r^{\eta}_{t,0,i}
    &=z_{t,i},\\
  g^\eta_{t,l,i}
    &=\nabla_s E_\eta\!\left(
      r^\eta_{t,l-1,i};o_t,B^\eta_{t-1},a_{t-1}
    \right),\\
  r^\eta_{t,l,i}
    &=r^\eta_{t,l-1,i}
      +\alpha_l g^\eta_{t,l,i}
      +\sqrt{2\alpha_l}\,\epsilon_{t,l,i},
      \qquad l=1,\ldots,L,\\
  s^\eta_{t,i}
    &=r^\eta_{t,L,i},\qquad
  q^\eta_{t,i}
    =\nabla_s E_\eta\!\left(
      s^\eta_{t,i};o_t,B^\eta_{t-1},a_{t-1}
    \right),\\
  B^\eta_t
    &=\{(s^\eta_{t,i},q^\eta_{t,i})\}_{i=1}^{K_b},
    \qquad c^\eta_t=E_\phi(B^\eta_t)
  \label{eq:belief-replay}
\end{align}
を順に再計算する。全ULA stepと最後のscore計算で
\verb|create_graph=True|を使うため、PPO/value lossから粒子生成経路、
初期energy、観測energy、遷移energyへpathwise gradientが流れる。noiseの分布はparametersに依存しないので、
保存noise自体は微分しない。

PPOが固定dataとして扱うのは$o_t$、実現済み$a_{t-1}$、$r_t$、終了mask、
old log-prob、advantageとreturnである。環境状態遷移や観測生成を再現するモデルは使わず、
環境へは微分しない。この意味で学習はmodel-freeでありながら、agent内部の
belief recurrenceには全勾配を流す。

\subsection{\texttt{full}と\texttt{tbptt\_1}}

\texttt{full}では、rollout開始時にepisodeが継続中なら、そのepisodeの開始から
rollout境界までの観測、実現済み行動、noiseからなるprefixも保持する。
各PPO minibatchでepisode-startから対象loss時刻まで式\eqref{eq:belief-replay}を
現在parametersで再生し、真のepisode resetだけでgraphを切る。
したがってrandomにflattenした環境時刻を独立sampleとして扱ってはならず、
時間順のsequence minibatchを用いる。

\texttt{tbptt\_1}は同じforward再生を用いるが、各環境時刻の冒頭だけ
\begin{equation}
  B^{\mathrm{in}}_{t-1}=\stopgrad(B^\eta_{t-1})
  \label{eq:tbptt-one}
\end{equation}
とする。現在時刻内の$L$個のULA stepと最後のscore計算はdetachしない。
よって両modeは同じnoiseとparametersに対して同一の粒子・action distributionを返し、
違いは時間方向のgradient範囲だけである。

optimizer stepの後には、それまでの未終了episode prefixを更新後parametersで再生し、
次rolloutのbelief contextを再構築する。これにより、更新前parametersで作った
stale particleを次rolloutへ持ち越さない。PPOの複数epochでも計算graphを使い回さず、
各optimizer stepのcurrent parametersでsequenceを再生する。

\texttt{full}のgraph memoryはepisode長に比例するため、実装では
non-reentrant activation checkpoint
\verb|torch.utils.checkpoint(..., use_reentrant=False)|を用いる。
backward時にforward activationを再計算するだけなので、detachやno-gradによる近似と異なり、
数学的な勾配は変えない。計算量は増えるためwall-clockとpeak GPU memoryをmode別に記録する。
さらにapproximate KL、clip fraction、global gradient normに加え、初期energy $g_0$、
観測energy $f$、遷移energy $u$それぞれのpre-clip gradient normを監視する。
初期時刻と再帰時刻を分けたscore/particleのL2 norm、各区分のparticle数、
score/particleのnon-finite数も記録する。これらは診断値であり、閾値によるoptimizer stepの
拒否には使わない。

productionでは論理PPO minibatchを$512=4\times128$遷移のまま保ち、
score beliefのsequence再生だけを、時間方向を切らずに$128$遷移ずつ4回処理する。
各memory microbatchの平均lossを$1/4$倍してbackwardし、4回分のgradientを加算した後に
gradient clippingとoptimizer stepをそれぞれ1回だけ行う。したがって、浮動小数点の加算順による
微小差を除けば512遷移の平均lossに対する更新と同じであり、
\texttt{full}と\texttt{tbptt\_1}の時間方向のgradient境界も変えない。

GRUにも同じ時間範囲の二条件を実装する。\texttt{full}では、rollout開始時にepisodeが
継続中なら観測・直前行動からなるepisode prefixを保持し、current parametersで
episode開始から式\eqref{eq:gru-condition}を再生する。rollout境界ではgraphを切らず、
真のepisode resetだけを境界とする。\texttt{tbptt\_1}では同じ入力列とparametersを用いるが、
各環境時刻の冒頭で
\begin{equation}
  h^{\mathrm{in}}_{t-1}=\stopgrad(h_{t-1})
  \label{eq:gru-tbptt-one}
\end{equation}
とする。両modeのhidden、方策分布、valueというforward値は同一であり、差は時間方向の
gradient範囲だけである。GRUのPPO minibatchは固定長time chunkへ切らず、完全な時間軸を
保ったまま環境軸だけを分割する。optimizer step後は未終了episode prefixを更新後parametersで
再生し、次rollout用のonline hiddenを再構築する。

\section{学習アルゴリズム}

\begin{algorithm}[H]
\small
\caption{Belief-conditioned DDPM policyのDPPO型from-scratch学習}
\label{alg:training}
\begin{algorithmic}[1]
\Require 並列環境、rollout長$T$、粒子数$K_b$、ULA step数$L$、拡散step数$N$、gradient mode
\Ensure 更新済み信念energy、set encoder、denoiser、critic
\State 全networkをランダム初期化し、episode初期粒子を設定する
\For{各PPO iteration}
  \State rollout用parametersをold policyとして固定する
  \For{$t=0,\ldots,T-1$}
    \State $z_{t,1:K_b},\epsilon_{t,1:L,1:K_b}$を生成・保存してULAを$L$ step実行する
    \State $q_{t,i}=\nabla_sE_{\eta_{\old}}(s;\cdot)|_{s=\bar s_{t,i}}$と$c_t$を計算する
    \State $x_t^N\sim\Normal(0,I)$を生成する
    \For{$n=N,N-1,\ldots,1$}
      \State $x_t^{n-1}\sim p_{\theta_{\old}}(\cdot\mid x_t^n,c_t)$
      \State $(x_t^n,x_t^{n-1},\ell^{\old}_{t,n})$を保存する
    \EndFor
    \State $a_t=T(x_t^0)$を環境へ渡し、$o_t,a_t,r_t,d_t,V_{\psi_{\old}}(c_t)$を保存する
  \EndFor
  \State 式\eqref{eq:gae}で$\widehat A_t^{\env},\widehat R_t$を計算し、advantageを正規化する
  \For{各PPO epochと論理sequence minibatch}
    \State gradientを初期化する
    \For{4個のmemory microbatch（各128遷移）}
      \State episode-startから保存noiseで式\eqref{eq:belief-replay}をcurrent parametersで再生する
      \If{gradient modeが\texttt{tbptt\_1}}
        \State 各環境時刻の入力beliefだけ式\eqref{eq:tbptt-one}でdetachする
      \EndIf
      \State 全保存遷移のnew log-probと式\eqref{eq:step-ratio}の$\rho_{t,n}$を計算する
      \State 式\eqref{eq:ppo-total}のlossを$1/4$倍し、checkpoint付きbackpropagationで加算する
    \EndFor
    \State 加算gradientを1回clipし、optimizer stepを1回行う
  \EndFor
  \State 未終了episode prefixをcurrent parametersで再生し、次rolloutのcontextを再構築する
\EndFor
\end{algorithmic}
\end{algorithm}

\section{行動選択：argmaxを使わない}

式\eqref{eq:intractable-marginal}の密度を直接評価できないため、通常の拡散方策では
\begin{equation}
  \argmax_a\pi_\theta(a\mid B_t)
  \label{eq:invalid-argmax}
\end{equation}
を計算できない。学習中は必ず逆拡散chainをsampleして
$a_t\sim\pi_\theta(\cdot\mid B_t)$とする。

評価時の再現可能な決定論的規約として、本実装は
\begin{equation}
  x_t^N=0,\qquad
  x_t^{n-1}=\mu_\theta(x_t^n,n,c_t),\qquad
  a_t=T(x_t^0)
  \label{eq:mean-chain}
\end{equation}
という\textbf{mean-chain}を使う。これは各conditional Gaussianの平均を順に通す規約であり、
周辺分布の平均やMAPと等しいとは限らない。報告時には
stochastic evaluationとmean-chain evaluationを区別する。
さらにscore系では、mean-chainにしてもbeliefの初期粒子とULA noiseはsampleされる。
したがって「決定論的」はaction readoutだけを指し、agent全体の軌跡が確定するという意味ではない。
\texttt{mean\_chain}はreturn等の性能指標ではなく、このaction readout規約の名前である。
tanh Gaussianを使う非active ablationの\texttt{mean\_action}も同様に、pre-tanh Gaussian平均を
tanhとaction scaleへ通すreadout名であり、metricではない。activeな連続比較は
score TransformerとGRUの双方が\texttt{mean\_chain}を使うため、deterministic readoutも揃う。
それでも事前に一つのprimary metricを選ぶなら、探索を含む学習方策そのものを評価するため、
held-out episodeにおけるstochastic evaluationの最終returnをtraining seed間で集約した値を用いる。

\section{標準DPPOとの差とlimitations}

\subsection{from-scratch設定}

RenらのDPPO\cite{ren2025dppo}の中心設定は、デモンストレーションでbehavior cloningした
Diffusion PolicyをRLでfine-tuneするものである。早期denoising stepを凍結し、
最後の$N'<N$ stepだけ更新する設計も事前学習を前提とする。
本実験では、信念energy、encoder、denoiser、criticをすべてランダム初期化し、
全denoising stepを環境報酬だけから学習する。DPPOの付録ではfrom-scratch実験もあるが、
Gaussian PPOより学習が遅く、10-step diffusionでwall-clockが約6倍と報告されている。
したがって本手法は「DPPOそのもの」ではなく、
\textbf{DPPO型のper-denoising-step PPOをfrom scratchで用いる手法}と記述する。

\subsection{主な制約}

\begin{enumerate}
  \item \textbf{sample efficiency:}
    環境1 stepを$N$個の内側stepへ展開するため、Gaussian PPOより計算量が大きい。
    デモなしの複雑なtaskでは探索が成立しない可能性がある。
  \item \textbf{recurrent gradientの計算量:}
    \texttt{full}はepisode-startから全ULA stepを再計算するため、時間、粒子数、ULA step数に対して
    高コストである。\texttt{tbptt\_1}は時間方向を切るablationであり、full-gradientの代替結果として
    混同せずmode別に報告する。
  \item \textbf{beliefの非同定性:}
    PPO報酬だけでは$f_\eta,u_\eta$が真の観測尤度・遷移密度になる保証はない。
    得られるものは制御に有用なtask-oriented latent beliefである。
  \item \textbf{variance floorと有界化:}
    $\sigma_{\min}$が小さすぎるとratioが不安定になり、大きすぎると最終行動が過度にランダムになる。
    $\tanh$の飽和率も記録する必要がある。
  \item \textbf{周辺尤度とMAPの欠如:}
    per-step PPOは拡張MDP上のsurrogateである。最終actionの周辺KLやMAPを直接制御しない。
  \item \textbf{denoising index:}
    保存chainを$(x^N,x^{N-1},\ldots,x^0)$の順に並べる場合、
    advantage重みは$(\gamma_{\den}^{N-1},\ldots,\gamma_{\den}^{0})$の順でなければならない。
    逆順にするとnoisy stepを強調してしまう。
  \item \textbf{比較の必要性:}
    実装上は、同じscore Transformerを持つ\texttt{score\_gaussian}に加え、
    observation MLPとoracle-stateのtanh Gaussian PPOも同じenvironment-step budgetで比較できる。
    ただしactive productionは\texttt{score\_transformer}とGRUに限定し、他の実装は
    local smokeで経路だけを確認する。activeな連続比較ではGRUにも提案法と同じDiffusionPolicy、
    per-denoising-step PPO、schedule、advantage重み、\texttt{mean\_chain} readoutを用いる。
    productionのparameter数も最大0.116\%差へ合わせるが、各手法を個別に十分tuningした比較ではない。
  \item \textbf{action dimension:}
    PendulumとMountainCarの1次元\texttt{Box}に加え、5次元Light-Darkを比較する。
    ただし無限次元行動や多峰なaction distributionに対するdiffusionの利点までは検証しない。
  \item \textbf{Light-Darkの報酬設計:}
    Light-Darkの報酬は遷移\emph{前}の状態と行動に課金する二次コストと、終了stepの
    goal bonusからなる:
    \[
      r_t=-\left(c_x\lVert x_t\rVert^2+c_u\lVert u_t\rVert^2\right)
        +B\,\mathbf 1\!\left\{\lVert x_{t+1}\rVert\leq r_g\right\},
    \]
    既定は$c_x=c_u=0.5$、$r_g=0.25$、$B=0$である。力学$x'=x+u$は決定論的なので事後平均は
    行動履歴の和だけで追跡でき、観測の価値は不確実性の縮小、すなわちgoal到達確率にしか現れない。
    $B=0$ではその到達確率へ報酬が付かないため、production条件（$N=5$、light at $x[0]=5$、
    初期$\Normal(2\cdot\mathbf 1,0.5^2I)$、horizon 30）では光へ寄る情報収集が厳密劣後する。
    実測では dead reckoning の平均returnが$-31.8$（goal到達率$0.03$）、光で自己位置同定してから
    原点へ戻る方策が$-87.9$（到達率$0.90$）であった。
    したがって\textbf{既定パラメータのLight-Dark結果をbelief品質の証拠として読んではならない}。
    情報収集を最適化するには損益分岐の$65$--$85$を十分に上回る$B$（目安$100$以上）を
    \texttt{environment.tasks.light\_dark.goal\_bonus}へ設定するが、これは報酬定義そのものの
    変更であり、既定値で実行した既存campaignとは混在させられない。
\end{enumerate}

\section{実装と報告のチェック項目}

\begin{description}
  \item[拡散chain] $N$、$\beta$ schedule、$\sigma_{\min}$、$\gamma_{\den}$、単一行動かchunkか。
  \item[PPO] $\gamma_{\env}$、$\lambda$、clip幅、epoch、論理minibatch、memory microbatch、gradient accumulation、gradient clip、診断用KL。
  \item[信念] 粒子数、ULA step数・step size、particle/score clip、保存noise、gradient mode。
  \item[評価] action readoutのmean-chain/mean-action、共通指標のstochastic、episode数、seed、action saturation率。
  \item[計算量] 環境step数に加えてmode別wall-clock、peak GPU memory、denoising transition数。
  \item[比較] activeなscore Transformer/GRUの共通action headとparameter数、非activeなscore Gaussian、observation MLP、oracle-state、Deep Sets。
  \item[診断] $g_0/f/u$のpre-clip gradient norm、initial/recursive score・particle norm、sample数、non-finite数。
\end{description}

実装済み6手法のlocal smokeは、実行可能な23 method/task組とgradient mode 2種類からなる
46 cellsを1 seed・1 PPO updateで確認する。一方、active productionの比較行列は
\texttt{score\_transformer}とGRUの2手法$\times$4 tasks$\times$gradient mode 2種類
$\times$5 seedsの合計80 runsとし、集約表はmodeとmethod/taskの組ごとに16行を保持する。
active productionでscore beliefを持たないbaselineはGRUだけである。GRUも両mode labelで実行し、
\texttt{full}ではepisode prefixをrollout境界越しに再生し、\texttt{tbptt\_1}では各環境stepの
入力hiddenをdetachする。したがって両modeは同一forward・異なるtemporal gradient範囲を持つ、
score beliefと対応したGRU ablationである。

productionのGRUは入力encoder幅を$[348,185]$、hidden幅を176とする。
共通policy/value headを含むparameter数は次の通りで、最大絶対相対差は0.116\%である。
\begin{center}
\begin{tabular}{lrrr}
\toprule
task & score Transformer & GRU & 相対差 \\
\midrule
CartPole & 347,462 & 347,060 & $-0.116\%$ \\
Pendulum & 348,357 & 348,759 & $+0.115\%$ \\
MountainCarContinuous & 348,229 & 348,411 & $+0.052\%$ \\
Light-Dark ($N=5$) & 351,817 & 352,223 & $+0.115\%$ \\
\bottomrule
\end{tabular}
\end{center}
local設定は経路確認用の小型networkであり、このcapacity matchingの対象外である。

\begin{thebibliography}{9}
\bibitem{ren2025dppo}
Allen Z. Ren, Justin Lidard, Lars L. Ankile, Anthony Simeonov, Pulkit Agrawal,
Anirudha Majumdar, Benjamin Burchfiel, Hongkai Dai, Max Simchowitz,
``Diffusion Policy Policy Optimization,'' ICLR, 2025.
\url{https://proceedings.iclr.cc/paper_files/paper/2025/hash/c0749c39aaff9e9e4c91f7118bf21b1e-Abstract-Conference.html};
\url{https://arxiv.org/abs/2409.00588}.

\bibitem{schulman2017ppo}
John Schulman, Filip Wolski, Prafulla Dhariwal, Alec Radford, Oleg Klimov,
``Proximal Policy Optimization Algorithms,'' 2017.
\url{https://arxiv.org/abs/1707.06347}.

\bibitem{pardo2018timelimits}
Fabio Pardo, Arash Tavakoli, Vitaly Levdik, Petar Kormushev,
``Time Limits in Reinforcement Learning,'' ICML, 2018.
\url{https://proceedings.mlr.press/v80/pardo18a.html};
\url{https://arxiv.org/abs/1712.00378}.

\bibitem{schulman2016gae}
John Schulman, Philipp Moritz, Sergey Levine, Michael Jordan, Pieter Abbeel,
``High-Dimensional Continuous Control Using Generalized Advantage Estimation,''
ICLR, 2016. \url{https://arxiv.org/abs/1506.02438}.

\bibitem{ho2020ddpm}
Jonathan Ho, Ajay Jain, Pieter Abbeel,
``Denoising Diffusion Probabilistic Models,'' NeurIPS, 2020.
\url{https://proceedings.neurips.cc/paper/2020/hash/4c5bcfec8584af0d967f1ab10179ca4b-Abstract.html}.

\bibitem{yang2025ncdpo}
Ningyuan Yang, Jiaxuan Gao, Feng Gao, Yi Wu, Chao Yu,
``Fine-tuning Diffusion Policies with Backpropagation Through Diffusion Timesteps,'' 2025.
\url{https://arxiv.org/abs/2505.10482}.
\end{thebibliography}

\end{document}
